% 第 18 章 Constraint Consistency
\bichapter{约束一致性}{Constraint Consistency}
\label{chap:cc}

约束一致性功能用于全芯片或模块级时序约束的分析与调试。基于规则检查的设计分析可识别缺失、无效、多余、低效或冲突的多种约束问题。约束一致性提供直观全面的图形界面（GUI），便于进一步调试。可使用大量内置规则运行分析，也可定义自定义一致性检查以扩展检查范围。

要了解约束一致性，请参阅：
\begin{itemize}
  \item 约束一致性概述
  \item 设计一致性检查
  \item 层次一致性检查
  \item 相关性一致性检查
  \item 附加分析特性
  \item 图形用户界面
  \item 教程
\end{itemize}

% ============================================================
\section{约束一致性概述}
\subsection*{Constraint Consistency Overview}

要了解约束一致性功能，请参阅：
\begin{itemize}
  \item 设计流程中的约束
  \item 调试特性
  \item 分析与调试方法学
  \item 启动约束一致性会话
  \item 支持的约束
  \item 报告约束接受情况
  \item 例外优先级顺序
  \item 分析违例
  \item 抑制违例
  \item \cmd{report\_constraint\_analysis} 命令
  \item 使用属性
  \item 自定义规则与报告
  \item 创建与使用用户定义规则
  \item 规则相关命令
  \item 读取不完整或不匹配网表的设计
\end{itemize}

\subsubsection{设计流程中的约束}
\subsection*{Constraints in the Design Flow}

正确完整的约束对逻辑综合、布局、布线等时序驱动实现步骤至关重要。约束描述工作环境及有效静态时序分析（独立或嵌入优化工具）所需信息。时序约束通常以 Synopsys Design Constraints（SDC）文件或更一般的 Tcl 脚本指定。错误或不完整约束可导致设计迭代浪费甚至硅片失效。应尽早发现并修复约束问题以降低风险、提高生产率。

约束一致性在门级分析约束，确定时钟、数据信号与常量在设计中的传播位置与方式。约束在各步骤常被修改（例如布局前约束经修改得到布局、CTS、布线、芯片组装约束）。每次创建或修改约束后应分析正确性与一致性。

约束一致性分析设计及其约束的正确性、完整性与一致性。读入并链接 Verilog 格式的顶层或模块设计，用与 Design Compiler 或 PrimeTime 相同的 SDC 命令施加约束。工具应用 SDC 约束，忽略约束脚本中的非 SDC 工具专用命令。读入、链接并约束后，可运行一条或多条约束一致性命令。\textbf{第一步}运行 \cmd{analyze\_design}，检查设计与约束并生成可在违例浏览器中查看的摘要报告。

从摘要列表可点击违例查看详情；从详细报告可点击链接定位 SDC 脚本中的约束命令、查看原理图相关区域或获取修复建议。检测的违例类型包括：
\begin{itemize}
  \item 错误电容值、错误指定的时钟/生成时钟
  \item 错误的时钟 latency 与转换时间
  \item 错误或未指定的驱动单元
  \item 错误指定的时序例外、输入/输出延迟
  \item 网表问题（如冲突驱动器）、冲突 case analysis 值
\end{itemize}

进一步诊断可用：\cmd{analyze\_paths}、\cmd{analyze\_clock\_networks}、\cmd{analyze\_unclocked\_pins}、\cmd{report\_case\_details}；\cmd{report\_analysis\_coverage}、\cmd{report\_exceptions}、\cmd{compare\_block\_to\_top} 有助于调试多场景约束、多重时序例外与层次设计中的冲突约束。

\subsubsection{调试特性}
\subsection*{Debugging Features}

重要特性包括：覆盖广泛潜在约束问题的规则集；用户定义规则；单进程多场景分析；数百万实例级设计的高效容量与性能；强大诊断与详细分析命令；丰富的集合命令与对象属性；符合 SDC 2.0 的规则检查器；支持循环/变量/条件的交互式 Tcl shell；含源文件名与行号的 GUI 输出；与 DC/ICC 等实现工具在约束解释、时钟/信号/常量传播及界面方面的高度一致。

\subsubsection{分析与调试方法学}
\subsection*{Analysis and Debugging Methodology}

\textbf{输入：}Verilog 网表、Synopsys \texttt{.db} 库、SDC 或 PrimeTime Tcl 约束、Tcl 控制脚本。

\textbf{输出：}文本报告；以及可用于交互式调试、自定义报告或条件脚本的查询命令与对象属性。

典型流程：
\begin{enumerate}
  \item 启动：\texttt{pt\_shell -constraints}（见“启动会话”）
  \item 读设计：\cmd{set\_app\_var search\_path}、\cmd{link\_path}、\cmd{read\_verilog}
  \item 链接：\cmd{current\_design}、\cmd{link\_design}
  \item 各场景施加约束：\cmd{source}（Tcl）或 \cmd{read\_sdc}（SDC）
  \item 分析：\cmd{analyze\_design}
  \item 诊断：专用命令（见“附加分析特性”）
  \item 按需重复分析与诊断
\end{enumerate}

\figplaceholder{Figure 229: Typical Constraint Consistency Flow}{典型约束一致性流程}{fig:cc-flow}

% ---- Starting session ----
\subsubsection{启动约束一致性会话}
\subsection*{Starting a Constraint Consistency Session}

约束一致性为 GUI 驱动，在 Linux 下运行。子主题：启动会话、许可证、Setup 文件、命令日志文件。

\paragraph{启动会话} 使用 \cmd{pt\_shell -constraints}。自动检出 PrimeTime SI 许可证，启动 GUI，出现 \texttt{ptc\_shell>} 提示符。若启动失败，检查：安装正确、PrimeTime 在 PATH 中、许可证服务器运行、许可证文件有效且含 PrimeTime SI 许可证。结束会话：\cmd{quit} 或 \cmd{exit}。

\paragraph{许可证} 启动 \texttt{ptc\_shell} 需 PrimeTime SI 许可证；会话退出自动检入。另需同一服务器上有 PrimeTime 许可证（不检出，但必须存在）。

\paragraph{Setup 文件} 每次启动执行 \texttt{.synopsys\_gca.setup}，按序检查：Synopsys 安装 \texttt{admin/setup}、用户主目录、当前工作目录。用 \cmd{-no\_init} 可禁止执行 setup 文件。

\paragraph{命令日志} 会话历史写入 \texttt{ptc\_shell\_command.log}（覆盖同名文件）。可在 setup 中用 \cmd{sh\_command\_log\_file} 指定其他名称（会话中不可改）。

\subsubsection{支持的约束}
\subsection*{Supported Constraints}

支持 SDC 与 Tcl 格式全部约束。Tcl 约束用 \cmd{source} 加载；SDC 用 \cmd{read\_sdc}。支持多场景：先 \cmd{create\_scenario}，再为各场景加载约束。

示例（Tcl 片段）：
\begin{lstlisting}
foreach_in_collection i [all_clocks] {
  echo [format "Clock : '%s'" [get_attribute $i full_name]]
  set_clock_uncertainty 1.5 $i
}
\end{lstlisting}

多场景加载示例：
\begin{lstlisting}
set all_scenarios {system myscenario_1 myscenario_2}
foreach scenario $all_scenarios {
  create_scenario ${scenario}
  source ./scripts/common.tcl
  source ./scripts/${scenario}_constraints.tcl
}
\end{lstlisting}

\subsubsection{报告约束接受情况}
\subsection*{Reporting Constraint Acceptance}

\cmd{report\_constraint\_analysis -include statistics} 报告分析接受的约束，快速确认约束是否正确加载。输出含各约束类型的 Accepted/Rejected/Total 计数。

\subsubsection{例外优先级顺序}
\subsection*{Exception Order of Precedence}

冲突时，约束一致性\textbf{按路径}（非按命令）独立应用优先级规则。例如 \cmd{set\_max\_delay -from A} 与 \cmd{set\_false\_path -to B}：对起自 A 且止于 B 的路径，\cmd{set\_false\_path} 优先；仅起自 A 的路径仍受 \cmd{set\_max\_delay} 约束。

\textbf{例外类型优先级}（高到低）：
\begin{enumerate}
  \item \cmd{set\_false\_path}
  \item \cmd{set\_max\_delay} / \cmd{set\_min\_delay}
  \item \cmd{set\_multicycle\_path}
\end{enumerate}

不视为冲突可同时有效的组合：两个 \cmd{set\_false\_path}；\cmd{set\_min\_delay} 与 \cmd{set\_max\_delay}；\cmd{set\_multicycle\_path} 的 setup 与 hold。

\textbf{路径说明优先级}（高到低）：
\begin{enumerate}
  \item \cmd{-from}/\cmd{-rise\_from}/\cmd{-fall\_from} pin
  \item \cmd{-to}/\cmd{-rise\_to}/\cmd{-fall\_to} pin
  \item \cmd{-through}/\cmd{-rise\_through}/\cmd{-fall\_through}
  \item \cmd{-from}/\cmd{-rise\_from}/\cmd{-fall\_from} clock
  \item \cmd{-to}/\cmd{-rise\_to}/\cmd{-fall\_to} clock
\end{enumerate}

组合示例（高到低）：\texttt{-from pin -to pin} $>$ \texttt{-from pin -to clock} $>$ \texttt{-from pin} $>$ \texttt{-from clock -to pin} $>$ \texttt{-to pin} $>$ \texttt{-from clock -to clock} $>$ \texttt{-from clock} $>$ \texttt{-to clock}。

用 \cmd{report\_exceptions -ignored} 列出被忽略的时序例外。

% ---- Analyzing violation + Suppressing ----
\subsubsection{分析违例}
\subsection*{Analyzing a Violation}

在违例浏览器中选择违例，在信息窗格查看详情。\figplaceholder{Figure 230--234}{违例浏览器、信息窗格链接、原理图与引脚标注}{fig:cc-violation}

可用 Debugging Help、Fix Suggestion 链接或 SDC 源文件链接进一步调试；查阅在线帮助中的规则参考。

\subsubsection{抑制违例}
\subsection*{Suppressing Violations}

可完全\textbf{禁用规则}，或对特定实例\textbf{豁免（waiver）}违例。GUI 与 \cmd{create\_waiver}、\cmd{report\_waiver}、\cmd{remove\_waiver}、\cmd{write\_waiver} 支持创建、报告、删除与写出豁免。

\figplaceholder{Figure 235--245}{违例抑制流程、豁免配置、已豁免违例图标}{fig:cc-waiver-flow}

\textbf{抑制流程：}可在 \cmd{analyze\_design} 前后抑制。分析后抑制的违例在浏览器中带特殊图标；分析前完全禁用的规则不出现在报告中；实例豁免仍报告但带图标。豁免后浏览器状态立即更新。

\textbf{禁用规则：}违例浏览器右键 “Disable rule”，或 Design $>$ Waiver Configuration 取消勾选，或 \cmd{disable\_rule}。

\textbf{豁免特定违例：}浏览器右键 “Waive violation”，或 Waiver Configuration 中 Create waiver，或 \cmd{create\_waiver}。\cmd{hide\_waived\_violations false} 时 View 菜单可切换隐藏已豁免违例。

\paragraph{\cmd{create\_waiver} 命令} 可指定多个 \cmd{-condition} 与 \cmd{-not\_condition}；只有全部条件均满足时才抑制违例。每个选项的实参必须是二元素列表：规则参数名，以及网表对象或约束的集合。

示例 1——豁免 CLK* 时钟上的 CLK\_0026：
\begin{lstlisting}
ptc_shell> create_waiver -rule CLK_0026 \
  -condition [list "clock" [get_clocks CLK*]]
\end{lstlisting}

示例 2——豁免除 CLK3 外所有时钟：
\begin{lstlisting}
ptc_shell> create_waiver -rule CLK_0026 \
  -not_condition [list "clock" [get_clocks CLK3]]
\end{lstlisting}

示例 3——多参数（时钟 CLK1/CLK2，引脚非 U1/A、U1/B）：
\begin{lstlisting}
ptc_shell> create_waiver -rule CLK_0003 \
  -condition [list "clock" [get_clocks {CLK1 CLK2}]] \
  -not_condition [list "pin" [get_pins U1/A] [get_pins U1/B]]
\end{lstlisting}

示例 5——例外参数：
\begin{lstlisting}
ptc_shell> create_waiver -rule EXC_0001 \
  -condition [list "exception" [get_exceptions -from U1/A -to U2/Z]] \
  -condition [list "rise_or_fall" "rise"]
\end{lstlisting}

用 \cmd{report\_rule -parameters rule\_name} 查询规则参数。B2T 与 S2S 规则豁免需指定 \cmd{-design1}/\cmd{-scenario1}/\cmd{-design2}/\cmd{-scenario2} 及 top/block 对象上下文。

\textbf{实例级豁免：}层次浏览器右键 Waive Instance，或 Waiver Configuration $>$ Instance Waivers；命令行用 \cmd{-cells}：
\begin{lstlisting}
ptc_shell> create_waiver -name my_waiver_1 -cells [get_cell U1/U1]
\end{lstlisting}

用法：实例级豁免应用于当前场景；\cmd{-all\_scenarios} 应用于设计所有场景。仅当违例相关对象完全包含在指定实例内才豁免；部分通用规则与用户定义规则不支持实例级豁免。

\cmd{write\_waiver -output file [-force]} 将豁免写入 Tcl 供后续 \cmd{source} 恢复。

% ---- report_constraint_analysis + Attributes ----
\subsubsection{\cmd{report\_constraint\_analysis} 命令}
\subsection*{report\_constraint\_analysis Command}

生成汇总设计中违例、用户消息与规则信息的文本报告；按分析类型、设计、场景排序。

\begin{itemize}
  \item \cmd{-include}：\texttt{violations}、\texttt{statistics}、\texttt{user\_messages}、\texttt{rule\_info} 等
  \item \cmd{-style full|summary}
  \item \cmd{-rules}/\cmd{-rule\_types} 过滤
  \item \cmd{-format csv} 输出 CSV
\end{itemize}

多设计层次比较结果在报告中以 Top/Block 场景分组显示 B2T 违例。

\subsubsection{使用属性}
\subsection*{Using Attributes}

属性是与设计对象关联的字符串或值。\textbf{命令：}\cmd{list\_attributes}、\cmd{get\_attribute}、\cmd{report\_attribute}。\cmd{get\_attribute} 限于单对象单集合。

属性组（表 43）包括：\texttt{cell}、\texttt{clock}、\texttt{clock\_group}、\texttt{design}、\texttt{exception}、\texttt{pin}、\texttt{port}、\texttt{net}、\texttt{rule}、\texttt{rule\_violation}、\texttt{scenario}、\texttt{timing\_arc}、\texttt{lib\_*} 等。默认只读。

\cmd{get\_timing\_arcs} 创建时序弧集合供自定义报告；\cmd{get\_lib\_timing\_arcs} 创建库时序弧集合。示例——查找 U1 上被禁用的 positive\_unate 弧：
\begin{lstlisting}
ptc_shell> set arcs [get_timing_arcs -of_objects U1 \
  -filter "sense == positive_unate"]
ptc_shell> foreach_in_collection arc $arcs {
  echo [get_attribute $arc is_disabled]
}
\end{lstlisting}

% ---- Custom rules (condensed) ----
\subsubsection{自定义规则与报告}
\subsection*{Customizing Rules and Reports}

可用 \cmd{define\_rule}、\cmd{define\_ruleset} 扩展检查范围；\cmd{report\_rule}、\cmd{enable\_rule}/\cmd{disable\_rule}、\cmd{set\_rule\_property} 管理规则（含容差 \texttt{tolerance}）。\cmd{analyze\_design -rules} 可限定检查的规则子集。

\subsubsection{创建与使用用户定义规则}
\subsection*{Creating and Using User-Defined Rules}

用户定义规则用 Tcl 过程检测自定义条件；可访问设计与约束对象属性。定义后通过 \cmd{analyze\_design} 或专用规则集运行。详见在线帮助 “User-Defined Rules”。

\subsubsection{规则相关命令}
\subsection*{Rule-Related Commands}

\begin{itemize}
  \item \cmd{report\_rule [pattern]} — 规则状态与描述
  \item \cmd{enable\_rule}/\cmd{disable\_rule} — 启用/禁用
  \item \cmd{set\_rule\_property -tolerance} — B2T/S2S 比较容差（默认 1e-5\%）
  \item \cmd{analyze\_design}、\cmd{compare\_block\_to\_top}、\cmd{compare\_constraints} — 各检查不同规则集
\end{itemize}

\subsubsection{读取不完整或不匹配网表的设计}
\subsection*{Reading Designs With Incomplete or Mismatched Netlists}

\cmd{read\_verilog} 与 \cmd{link\_design} 可处理部分不匹配；\cmd{report\_design\_mismatch} 报告主实例与引用不一致（如 DMM-038 引脚不存在、DMM-905 库中找不到单元）。应在分析前尽量解决链接错误。

% ============================================================
\section{设计一致性检查}
\subsection*{Design Consistency Checking}

主题：规则与违例、内置规则。

\subsubsection{规则与违例}
\subsection*{Rules and Violations}

对已加载、已链接、已约束的设计检查设计与约束相关条件。\textbf{规则（rule）}为检查项；\textbf{违例（violation）}为未满足条件的一次出现。违例严重级别：Information、Warning、Error。规则代码为三到四字母+下划线+四位数字（如 \texttt{CLK\_0003}：生成时钟主源引脚无时钟）。

\cmd{analyze\_design} 检查标准规则集；\cmd{compare\_block\_to\_top} 与 \cmd{compare\_constraints} 各检查不同规则集。可 \cmd{disable\_rule} 禁用规则；部分规则默认禁用以节省运行时间。\textbf{豁免（waiver）}在特定条件下跳过检查。除内置规则外可定义用户规则。

\subsubsection{内置规则}
\subsection*{Built-In Rules}

启动时自动加载表 44 所示内置规则：

\begin{table}[htbp]
\centering
\caption{内置规则分类（表 44）}
\small
\begin{tabular}{|l|l|}
\hline
边界条件 & CAP\_xxxx 电容；DRV\_xxxx 驱动；EXD\_xxxx 外部延迟 \\
约束/例外分析 & CAS\_xxxx case；EXC\_xxxx 时序例外 \\
时钟 & CGR\_xxxx 组；CLK\_xxxx 属性；CNL\_xxxx 网络 latency；\\
& CTR\_xxxx 转换；CSL\_xxxx 源 latency；UNC\_xxxx uncertainty \\
通用 & CMP\_xxxx 工具兼容；DES\_xxxx 设计约束；HIER\_xxxx 层次；\\
& LOOP\_xxxx 环；PRF\_xxxx 性能；NTL\_xxxx 网表；UNT\_xxxx 库单位 \\
\hline
\end{tabular}
\end{table}

详细规则描述见 Constraint Consistency 在线帮助。示例：\texttt{CLK\_0003} 表示生成时钟因主源引脚无时钟而无法展开。

% ============================================================
\section{层次一致性检查}
\subsection*{Hierarchical Consistency Checking}

设计团队采用多种策略创建模块级与顶层约束；自底向上流程中模块约束可传播到顶层，迭代中顶层与模块均可能手工插入新约束。\cmd{compare\_block\_to\_top} 比较模块级与顶层约束是否匹配。

各模块带有综合/布局所用约束（时钟、例外、边界条件）。层次一致性对\textbf{同一网表}加载两套约束并比较，生成 B2T 规则违例，可在违例浏览器中调试；单会话可保留多块 B2T 结果。

\figplaceholder{Figure 246--250}{多块违例浏览器、B2T 原理图与双 SDC 查看器}{fig:cc-b2t}

\subsubsection{比较模块级与顶层约束}
\subsection*{Comparing Block- and Top-Level Constraints}

复用 IP 或块时须保证并入更大设计后块约束仍有效。\cmd{compare\_block\_to\_top} 找层次约束不一致；\cmd{analyze\_design} 检查单套约束内部问题。

\textbf{方法学：}
\begin{itemize}
  \item \textbf{保持多个链接设计：}默认 \cmd{link\_design} 会移除先前链接设计；B2T 流程需同时保留块设计（块约束）与顶层设计（顶约束）。用 \cmd{link\_design -add}；\cmd{current\_design} 切换。
  \item 流程脚本示例（Example 56）：
\end{itemize}

\begin{lstlisting}
set_app_var search_path [list . ${TOP_DIR}/libs ${TOP_DIR}/design]
set_app_var link_path [list * libs.db]
read_verilog ${TOP_DIR}/design/ChipLevel_no2blocks.v
read_verilog ${TOP_DIR}/design/Multiply16x16.v
link_design ChipLevel
link_design -add Multiply16x16
current_design ChipLevel
source ${TOP_DIR}/dc_outdir/ChipLevel_propagate.sdc -echo
current_design Multiply16x16
source ${TOP_DIR}/design/Multiply16x16.sdc -echo
current_design ChipLevel
compare_block_to_top -block_design [get_designs Multiply16x16]
\end{lstlisting}

\textbf{GUI 步骤：}加载设计 $\rightarrow$ 顶层与块级分别 \cmd{analyze\_design} 确保干净 $\rightarrow$ \cmd{compare\_block\_to\_top} $\rightarrow$ Analysis $>$ B2T Mismatch Browser $\rightarrow$ 调试（可豁免 B2T 违例）。

支持：多场景比较、同一块的多个实例、创建豁免、文本报告（\cmd{report\_constraint\_analysis}）。

\subsubsection{内置层次一致性规则}
\subsection*{Built-In Hierarchical Consistency Checking Rules}

B2T 规则前缀 \texttt{B2T\_}，分类包括：B2T\_CLK\_xxxx 时钟；B2T\_EXC\_xxxx 例外；B2T\_CAS\_xxxx case；B2T\_EXD\_xxxx 外部延迟；B2T\_UNC\_xxxx uncertainty；B2T\_DIS\_xxxx 禁用对象/弧；B2T\_OPC\_xxxx 工作条件等。比较规则可设容差属性；默认 1e-5\%，用 \cmd{set\_rule\_property} 修改。

% ============================================================
\section{相关性一致性检查}
\subsection*{Correlation Consistency Checking}

设计开发中约束文件可能被修改，导致原版与修改版不一致。\textbf{相关性一致性}用同一网表比较两套约束，或比较两个设计及各自约束集。

\subsubsection{比较两套设计约束}
\subsection*{Comparing Two Sets of Design Constraints}

用 \cmd{compare\_constraints} 比较。将两套约束建模为\textbf{场景}，在当前设计上创建后比较。

\cmd{compare\_constraints} 检查：\cmd{create\_clock}/\cmd{create\_generated\_clock}；\cmd{set\_input\_delay}/\cmd{set\_output\_delay}；\cmd{set\_false\_path}、\cmd{set\_multicycle\_path}、\cmd{set\_min\_delay}/\cmd{set\_max\_delay}、\cmd{set\_clock\_groups}；\cmd{set\_case\_analysis}；\cmd{set\_disable\_timing}；\cmd{set\_clock\_uncertainty}；\cmd{set\_clock\_transition}/\cmd{set\_input\_transition}；\cmd{set\_clock\_latency}；\cmd{set\_clock\_gating\_check}；\cmd{set\_sense}；\cmd{set\_load} 等。

\textbf{方法学：}
\begin{enumerate}
  \item \cmd{create\_scenario} 建模原版与修改版约束集
  \item \cmd{compare\_constraints -constraints1 ... -constraints2 ...}
  \item 用 S2S（SDC-to-SDC）浏览器识别违例；\cmd{report\_constraint\_analysis} 获取文本报告
  \item 按序修复：时钟 $\rightarrow$ \cmd{set\_disable\_timing} $\rightarrow$ \cmd{set\_case\_analysis} $\rightarrow$ 例外
  \item 重复比较
\end{enumerate}

\figplaceholder{Figure 252--254}{S2S 违例浏览器、双原理图与双场景 SDC}{fig:cc-s2s}

违例规则前缀 \texttt{S2S\_}。可禁用规则类型；禁用时 GUI 与文本报告均不显示。豁免见“抑制违例”。

\subsubsection{比较两个设计及两套约束}
\subsection*{Comparing Two Designs With Two Sets of Design Constraints}

链接两个功能等价设计并加载两套约束后，须先提供\textbf{名称映射文件（NMF）}再运行 \cmd{compare\_constraints}。可用 \cmd{-use\_clock\_map} 引用自定义时钟映射。需 PrimeTime-ELT 许可证。

\textbf{方法学：}
\begin{enumerate}
  \item 链接设计并加载约束（第二设计：\cmd{link\_design -add -latest design2}）
  \item \cmd{define\_name\_maps} 指定 NMF（仅 NMF 格式被识别）
  \item \cmd{compare\_constraints -constraints1 ... -constraints2 ... [-design2 design2]}
  \item \cmd{report\_constraint\_analysis} 与 GUI 调试
\end{enumerate}

NMF 语法：
\begin{lstlisting}
define_name_maps -application golden_sdc -design_name design2 \
  -columns {class pattern options names} \
  [list port IN [list] [list IN_2]] \
  [list cell BLK1 [list] [list BLK2]]
\end{lstlisting}

支持引脚、端口、网线、单元（含层次单元）映射。时钟映射示例：
\begin{lstlisting}
set_clock_map -clocks1 gclk1_hier -clocks2 gclk1_buf \
  -design2 design2 -scenario2 sdc2_mode1
\end{lstlisting}

% ============================================================
\section{附加分析特性}
\subsection*{Additional Analysis Features}

\subsubsection{\cmd{analyze\_paths} 命令}
\subsection*{analyze\_paths Command}

当 PrimeTime/DC/ICC 的 \cmd{report\_timing} 无法显示某些路径时，用 \cmd{analyze\_paths} 查明路径传播被阻断的原因及阻断点。可识别特定路径上的例外与约束，统计满足 from/through/to 条件的路径数量，并按受影响路径数对引脚排序。

\cmd{-path\_type full} 且 GUI 打开时生成文本报告与原理图（\figplaceholder{Figure 255}{analyze\_paths 识别的路径}{fig:cc-analyze-paths}）。\cmd{-max\_endpoints} 限制原理图显示端点数。\cmd{-text\_only} 或 GUI 关闭时仅文本。

报告类型：
\begin{itemize}
  \item \textbf{默认：}到/经/自某对象集合的路径数与时钟交互（Example 64）
  \item \textbf{\cmd{-path\_type full}：}含起点/终点细分、逻辑级、途经点；若存在例外则显示例外名、脚本名与行号（Example 65--66）
  \item \textbf{\cmd{-path\_type summary}：}摘要（Example 67）
  \item \textbf{\cmd{-traverse\_disabled}：}追踪被禁用路径，显示数据传播阻断原因（Example 68）
\end{itemize}

示例：
\begin{lstlisting}
ptc_shell> analyze_paths -to Tranx/pay_count__reg[1]/D
ptc_shell> analyze_paths -to Tranx/pay_count__reg[1]/D -path_type full
ptc_shell> analyze_paths -to Tranx/CRC_tranx/pay_out_reg/D \
  -path_type full -traverse_disabled
\end{lstlisting}

\subsubsection{\cmd{analyze\_unclocked\_pins} 命令}
\subsection*{analyze\_unclocked\_pins Command}

分析引脚集合中缺失时钟定义或时钟传播被阻断的情况；默认分析设计中所有引脚。默认报告：分析摘要、可能缺失时钟定义的位置、时钟传播被阻断的位置。\cmd{-verbose} 列出无时钟的时钟引脚、case 禁用与寄存器禁用引脚。原理图 Pin annotation 菜单可开关时钟与 case 标注。

\subsubsection{\cmd{analyze\_clock\_networks} 命令}
\subsection*{analyze\_clock\_networks Command}

分析复杂时钟网络的传播，并生成报告，说明时钟网络从指定点出发或到达指定点时的传播情况以及约束所产生的影响。报告可覆盖实际时钟网络、使用 \cmd{-traverse\_disabled} 检测到的潜在时钟网络，以及生成时钟的源 latency 网络；这对于调试时钟定义以及由 \cmd{set\_disable\_timing} 或 case analysis 导致的时钟传播阻断尤其有用。

\subsubsection{\cmd{report\_case\_details} 命令}
\subsection*{report\_case\_details Command}

报告设计中 case analysis 传播结果：用户设置值、传播常量、冲突与未解析网络。用于调试 \cmd{set\_case\_analysis} 与 CAS\_xxxx 违例。

\subsubsection{\cmd{report\_clock\_crossing} 命令}
\subsection*{report\_clock\_crossing Command}

报告时钟域交叉路径统计：发射/捕获时钟对、路径数量、是否被 false path/clock groups 约束。用于验证跨时钟域约束完整性。

\subsubsection{\cmd{report\_analysis\_coverage} 命令}
\subsection*{report\_analysis\_coverage Command}

（PrimeTime 侧）报告分析覆盖率；约束一致性环境中用于检查约束导致的未测试路径。与 \cmd{check\_timing} 配合确保断言有效。

\subsubsection{\cmd{report\_exceptions} 命令（冗余约束）}
\subsection*{report\_exceptions Command (Redundant Constraints)}

识别冗余、被覆盖与主导的时序例外：
\begin{itemize}
  \item \cmd{report\_exceptions -ignored} — 被其他例外完全覆盖的例外
  \item \cmd{report\_exceptions -redundant} — 冗余例外
  \item \cmd{report\_exceptions -dominant} — 对至少一条路径起主导作用、移除会改变约束集的例外
  \item \cmd{report\_exceptions -ignored -verbose} — 显示主导忽略例外的详细信息
\end{itemize}

约束一致性不修改约束；编辑 SDC 后重新运行验证。应从 SDC 中删除被忽略的例外以节省后续会话运行时间与内存。

% ============================================================
\section{图形用户界面}
\subsection*{Graphical User Interface}

约束一致性 GUI 为查看与调试 PrimeTime、Design Compiler、IC Compiler 等工具约束问题的主要环境。

子节：使用 GUI、约束一致性窗口、用户消息浏览器、违例浏览器、B2T 违例浏览器、相关性一致性违例浏览器、豁免配置对话框、信息窗格、控制台、在线帮助、SDC 视图、层次浏览器与原理图视图、按名选择工具栏、属性对话框、原理图显示选项。

\subsubsection{使用 GUI}
\subsection*{Using GUI}

典型任务：启动 GUI $\rightarrow$ 读入/链接/约束/分析 $\rightarrow$ 深入分析违例路径。默认 \texttt{pt\_shell} 含 GUI；可用 \cmd{gui\_start}/\cmd{gui\_stop}。仅 shell：\texttt{pt\_shell -no\_gui}。View $>$ Palette/Toolbars 显示/隐藏界面元素。控制台含 \texttt{ptc\_shell} 提示符及 Log/History 标签。File $>$ Close GUI 关闭 GUI 保留终端会话；File $>$ Exit 退出。

\subsubsection{约束一致性窗口}
\subsection*{The Constraint Consistency Window}

窗口可平铺（Window $>$ Tile）、层叠（Cascade）、移动与调整大小。表 46 窗口类型：

\begin{table}[htbp]
\centering
\caption{约束一致性窗口类型（表 46）}
\small
\begin{tabular}{|L{3.2cm}|L{8.8cm}|}
\hline
\tblhead{窗口} & \tblhead{描述} \\
\hline
Violation Browser & 按场景与严重级别列出违例 \\
B2T Mismatch Browser & 顶层与块级约束差异 \\
S2S Mismatch Browser & 同一设计两套约束差异 \\
User Message Browser & 用户消息 \\
Information Pane & 当前选择的详细信息 \\
Console & 运行 \texttt{ptc\_shell}、查看日志与历史 \\
SDC View & 显示 SDC 并标记相关行 \\
Hierarchy/Schematic & 浏览层次、选择单元/网线/端口/引脚 \\
Path Schematic & 层次单元或路径的原理图 \\
Properties & 对象属性列表 \\
\hline
\end{tabular}
\end{table}

启动时默认停靠：违例浏览器、信息窗格、控制台（\figplaceholder{Figure 259}{启动时的约束一致性窗口}{fig:cc-gui-start}）。可按需打开层次浏览器（Window $>$ New Hierarchy Browser）。

工具栏提供常用菜单快捷方式；上下文变化时按钮启用/变灰。菜单栏：File、View、Selection、Highlight、Design、Analysis、Schematic、Window、Help（表 47）。

\subsubsection{违例浏览器}
\subsection*{Violation Browser}

主要调试工具；分层显示设计与约束违例及修复建议。支持网表相关与场景相关违例。显示顺序：Error、Warning、Information。默认每规则最多 1000 个违例标识节点；\cmd{display\_violations\_per\_rule\_limit} 可改。

多分析运行结果存于多个违例浏览器；Analysis Run Chooser（\figplaceholder{Figure 263}{分析运行选择器}{fig:cc-run-chooser}）选择查看哪次运行。首次打开时严重级别节点展开、违例标识节点折叠。

导航指南：单击查看信息窗格；双击展开规则下全部违例；使用过滤与搜索缩小列表；右键可禁用规则、豁免违例、查看 SDC/原理图。

\subsubsection{其他 GUI 组件（摘要）}
\subsection*{Other GUI Components}

\textbf{用户消息浏览器}（Analysis $>$ UserMessage Browser）：按严重级别排序；选择消息后在信息窗格显示 SDC 行号链接。

\textbf{B2T / S2S 浏览器：}与违例浏览器类似，提供并排显示的双原理图和双 SDC 视图，以及修复建议链接。

\textbf{豁免配置：}Design $>$ Waiver Configuration — 禁用规则、创建/编辑/删除豁免、实例豁免。

\textbf{信息窗格：}显示违例详情、超链接（SDC、原理图、Debugging Help、Fix Suggestion）、属性与例外摘要。

\textbf{SDC 视图：}点击链接打开 SDC 并高亮违例行。

\textbf{原理图：}Pin annotation 可显示 \texttt{case\_value}、\texttt{clock\_names}、\texttt{exceptions\_summary} 等；Attribute Group Manager 自定义显示属性组。

\textbf{在线帮助：}Help $>$ Online Help；\cmd{report\_rule CAS\_*} 等可快速查看规则摘要。

% ============================================================
\section{教程}
\subsection*{Tutorial}

本教程演示如何在教程设计上设置并运行约束一致性，并调查规则违例。假设读者已熟悉约束一致性关键组件。

子节：约束一致性概述、教程设计、启动、分析 ChipLevel 设计、调试约束问题、结束与重启会话。

\subsubsection{约束一致性概述（教程）}
\subsection*{Overview of Constraint Consistency}

违例浏览器（\figplaceholder{Figure 275}{约束一致性 GUI}{fig:cc-tut-gui}）为关键特性：\cmd{analyze\_design} 后按场景汇总 Error/Warning/Info；双击或 “+” 展开查看具体规则违例。信息窗格显示选中项详情与调查链接。控制台可交互输入命令并显示日志；History 标签列出已用命令。点击原理图链接可替换违例浏览器视图，用标签页切回。

\cmd{analyze\_design} 对照预定义规则检查已加载设计（时钟、case、例外、边界条件等）。表 50 规则类型与表 44 类似，另含 B2T\_xxx\_xxxx 块到顶规则。\cmd{report\_rule CAS\_*} 列出 case 规则摘要；详细描述见在线帮助。

\subsubsection{关于教程设计}
\subsection*{About the Tutorial Design}

教程设计 \textbf{ChipLevel} 实现多种二进制算术功能（加法、乘法等）。主要块：Adder16、CascadeMod、Comparator、Multiply8x8、Multiply16x16、MuxMod、PathSegment。设计处于前期开发阶段，时钟网络与约束仍在构建，适合用约束一致性尽早发现约束问题。

教程文件位于 PrimeTime 安装 \texttt{doc/gca/tutorial}：
\begin{lstlisting}
% cp -r $SYNOPSYS/doc/gca/tutorial .
\end{lstlisting}

目录结构：
\begin{verbatim}
tutorial/
  run_tutorial.tcl
  design/chiplevel.v, chiplevel.sdc
  libs/lsi_10k.db
\end{verbatim}

从 \texttt{tutorial} 目录运行 \texttt{run\_tutorial.tcl}。

\subsubsection{启动约束一致性}
\subsection*{Starting Constraint Consistency}

\begin{enumerate}
  \item 打开 Linux shell，进入 \texttt{user\_directory/tutorial}
  \item \texttt{pt\_shell -constraints}
\end{enumerate}

失败时检查安装、PATH、许可证服务器、PrimeTime 与 PrimeTime SI 许可证。

\subsubsection{分析 ChipLevel 设计}
\subsection*{Analyzing the ChipLevel Design}

步骤摘要：
\begin{enumerate}
  \item \cmd{set\_app\_var search\_path [list . ./libs ./design]}
  \item \cmd{set\_app\_var link\_path [list * lsi\_10k.db]}
  \item \cmd{read\_verilog chiplevel.v}
  \item \cmd{current\_design ChipLevel}
  \item \cmd{link\_design}
  \item \cmd{source ./design/chiplevel.sdc}（或 \cmd{read\_sdc}）
  \item \cmd{analyze\_design}
\end{enumerate}

\cmd{analyze\_design} 后违例浏览器显示结果。上述命令亦包含在 \texttt{run\_tutorial.tcl}（Example 84）中。

\subsubsection{调试约束问题}
\subsection*{Debugging Constraint Problems}

在违例浏览器中按严重级别展开 Error，选择具体违例（如 CLK\_xxxx、CAS\_xxxx）。使用信息窗格中的 SDC 链接、原理图链接、Debugging Help 与 Fix Suggestion。常见修复：补全时钟定义、修正 \cmd{set\_input\_delay}/\cmd{set\_output\_delay}、解决 case 冲突、删除无效 \cmd{set\_multicycle\_path} 对象（如 Q 引脚非有效端点）。

修改 \texttt{chiplevel.sdc} 后重新 \cmd{source} 并 \cmd{analyze\_design} 验证。可用 \cmd{report\_constraint\_analysis -include violations} 生成文本摘要。

\subsubsection{结束与重启教程会话}
\subsection*{Ending and Restarting the Tutorial Session}

\cmd{exit} 结束会话。保留 \texttt{ptc\_shell\_command.log} 可 \cmd{source} 重复会话。重新启动：\texttt{pt\_shell -constraints} 后 \cmd{source run\_tutorial.tcl}。

\begin{noteBox}
约束一致性（\texttt{pt\_shell -constraints}）与标准 PrimeTime STA 使用相同 SDC 语义，但专注于约束质量而非 signoff 时序。修复约束后应在 PrimeTime 中运行完整 STA 验证。
\end{noteBox}

\begin{seeAlsoBox}
SDC 格式见 Synopsys Design Constraints 应用说明；Tcl 见 \textit{Using Tcl With Synopsys Tools}；PrimeTime 报告与调试见第~\ref{chap:report}~章。
\end{seeAlsoBox}

